\section*{Hoofdstuk \ref{ch:b2:palaeoenvironments} --- Fossielen en paleomilieus}

\begin{solution}{exo:b2:palaeoenvironments:1}
(a) $112/135 = 0.830$ en $112/150 = 0.747$. (b) De dode schelpen op de bodem verzamelen zich over jaren
of langer: ze omvatten soorten die er vroeger leefden, of zelden, of die werden aangevoerd, dus de
\omterm{def:b2:palaeoenvironments:assemblages}{doodsassociatie} middelt vele generaties van de gemeenschap uit. (c) Groepen met een zacht lichaam, in
het bijzonder detritivoren en carnivoren zoals wormen: zonder skelet vergaan ze vóór de begraving.
\end{solution}

\begin{solution}{exo:b2:palaeoenvironments:2}
(a) $5 \times 165 = \qty{825}{kyr}$, ongeveer \qty{0.8}{Myr}, binnen de 0.7 tot 1 miljoen jaar.
(b) $183.73 - 182.77 = \qty{0.96}{Myr}$, en $0.96/0.165 = 5.8$: ongeveer zes horizonten. (c) Een
afstammingslijn verandert binnen veel minder dan een miljoen jaar; alleen horizonten scheiden de opeenvolgende
vormen, terwijl een zone ze samenvoegt.
\end{solution}

\begin{solution}{exo:b2:palaeoenvironments:3}
(a) \omterm{def:b2:palaeoenvironments:ichnofacies}{Nereites-ichnofacies}: de diepzeebodem van een turbidietsysteem, voorbij de platrand, waar
voedsel schaars is. (b) \omterm{def:b2:palaeoenvironments:ichnofacies}{Skolithos-ichnofacies}: vooroever of onderwateroever. (c) \omterm{def:b2:palaeoenvironments:ichnofacies}{Cruziana-ichnofacies}:
offshore, tussen de golfbasis bij mooi weer en de stormgolfbasis. (d) \omterm{def:b2:palaeoenvironments:ichnofacies}{Zoophycos-ichnofacies}: in het
Mesozoïcum het lagere plat of de helling, onder de stormgolfbasis.
\end{solution}

\begin{solution}{exo:b2:palaeoenvironments:4}
(a) Laag A bevat stenohaliene groepen: normaal zeewater. Laag B, arm aan soorten maar rijk aan
individuen, met plantenresten en geen van de stenohaliene groepen, wijst op brak water, bijvoorbeeld
een estuarium. (b) De ontbrekende groepen zijn talrijk, algemeen en goed bewaard in mariene lagen
van dezelfde ouderdom: hun gezamenlijke afwezigheid is een patroon, terwijl één fossiel bij toeval of
door slechte bewaring kan ontbreken. (c) Het sediment (getijde- of riviergeulen, moddermantels, wortels), de
\omterm{def:b2:palaeoenvironments:fossil-kinds}{sporenfossielen} (een \omterm{def:b2:palaeoenvironments:ichnofacies}{Cruziana-ichnofacies} met weinig soorten), en de zijdelingse overgang van laag B naar
mariene of naar rivierafzettingen.
\end{solution}

\begin{solution}{exo:b2:palaeoenvironments:5}
(a) $P = 1/2$ wanneer $kh/s = \ln 2$: $s = 0.05 \times 0.1/0.693 = \qty{0.0072}{m.yr^{-1}}$,
\qty{7.2}{mm} per jaar. (b) Een halvering van $s$ verdubbelt de exponent: $P = (1/2)^2 = 0.25$, de helft.
(c) $k$ wordt gehalveerd: $P = e^{-2.5} = 0.082$ in plaats van $e^{-5} = 0.0067$, $e^{2.5} = 12$ keer
waarschijnlijker. (d) Dikke, robuuste schelpen zijn oververtegenwoordigd, en des te meer naarmate de begraving trager is.
\end{solution}

\begin{solution}{exo:b2:palaeoenvironments:6}
(a) De bovengrens is $183.73 + 0.35 = \qty{184.08}{Ma}$: de \qty{184.2}{Ma} van de kaart ligt
erbuiten, met \qty{0.12}{Myr}. (b) $183.73 - 182.77 = \qty{0.96}{Myr}$. (c) $6/0.96 =
\qty{6.25}{m}$ per miljoen jaar. (d) Een biozone wordt gedefinieerd door fossielen, waarvan het verschijnen en
verdwijnen over grote gebieden bijna gelijktijdig is; de dikte van het gesteente dat in
die tijd werd afgezet, hangt af van de plaatselijke accumulatiesnelheid.
\end{solution}

\begin{solution}{exo:b2:palaeoenvironments:7}
(a) $0.49\,e^{-1046/3704} = 0.369$ en $0.49\,e^{-1166/3704} = 0.358$; gemiddelde $0.364$. (b) $120
\times (1 - 0.364) = \qty{76.3}{m}$. (c) $76.3/(1 - 0.49) = \qty{149.6}{m}$. (d) De factor
veronderstelt dat de hele herstelde laag de porositeit $\varphi_0$ van het oppervlak heeft; maar een laag
van \qty{150}{m} dik aan het oppervlak is aan haar basis minder poreus dan aan haar top, dus dezelfde korrels
vullen minder dan \qty{149.6}{m}.
\end{solution}

\begin{solution}{exo:b2:palaeoenvironments:8}
(a) Een volume-eenheid bevat $\varphi$ water en $1 - \varphi$ korrels: $\rho = \varphi
\rho_{\text{water}} + (1 - \varphi)\rho_{\text{korrel}}$, wat de formule geeft; $(2710 -
2573)/(2710 - 1000) = 0.080$. (b) $0.51\,e^{-956/4545} = 0.413$. (c) Cementatie: calciet dat
in de poriën neerslaat, wat carbonaten naast compactie ondergaan, neemt veel meer porositeit weg
dan begraving alleen.
\end{solution}

\begin{solution}{exo:b2:palaeoenvironments:9}
(a) Diepzeeturbidieten (\omterm{def:b2:palaeoenvironments:ichnofacies}{Nereites-ichnofacies}); een jurassische bodem van het lagere plat of de helling
(Zoophycos); een offshorebodem tussen de golfbases (Cruziana); een onderwateroever (Skolithos); een
moeras op het land. (b) Het water werd naar boven ondieper: de kust verschoof zeewaarts. (c) Een daling van het zeeniveau,
of een sedimentaanvoer die sneller is dan de daling van de bodem, zoals voor een groeiende delta; een zeespiegeldaling zou
op hetzelfde tijdstip in andere bekkens worden teruggevonden, een delta zou plaatselijk zijn en zijn sediment zou
van één kant komen. (d) In een concordante opeenvolging volgen naburige milieus elkaar op; een
sprong over de tussenliggende wijst op een tijdsleemte of op een breuk.
\end{solution}

\begin{solution}{exo:b2:palaeoenvironments:10}
(a) De tijd tot begraving is $h/s$ = 100, 10 en \qty{1}{yr}. Zacht lichaam: $e^{-100} \approx
\num{4e-44}$, $e^{-10} = \num{4.5e-5}$, $e^{-1} = 0.37$. Schelp ($k = \qty{0.01}{yr^{-1}}$):
$e^{-1} = 0.37$, $e^{-0.1} = 0.90$, $e^{-0.01} = 0.99$. (b) $s = kh/\ln 100 = 0.1/4.605 =
\qty{0.022}{m.yr^{-1}}$, \qty{22}{mm} per jaar. (c) Nauwelijks, over maanden; maar een stormlaag, een
turbidiet of een asregen begraaft de bodem in uren. (d) De constante sedimentatiesnelheid: de
tijd tot begraving is dan veel korter dan $h/s$ met de gemiddelde snelheid, en $P = e^{-kT}$, met $T$
de werkelijke tijd tot begraving, gaat naar één wanneer $T$ veel kleiner is dan $1/k$.
\end{solution}

\begin{solution}{exo:b2:palaeoenvironments:11}
(a) Een dunne laag $\mathrm dz$ op de diepte $z$ bevat de korrels $(1 - \varphi(z))\,\mathrm dz$;
aan het oppervlak vullen ze $\mathrm dz_0$ met $(1 - \varphi_0)\,\mathrm dz_0 = (1 -
\varphi(z))\,\mathrm dz$. (b) $\varphi = 0.63\,e^{-816/1960} = 0.415$; $(1 - 0.415)/0.37 =
1.581$, dus \qty{158.1}{m}. (c) \qty{9.5}{m} te veel: zoals in \cref{exo:b2:palaeoenvironments:7}
is de herstelde laag niet overal zo poreus als haar oppervlak. (d) De korrelsnelheid,
$h(1-\varphi)/\Delta t$, gebruikt alleen de huidige dikte en porositeit, die allebei gemeten zijn.
\end{solution}

\begin{solution}{exo:b2:palaeoenvironments:12}
(a) $9.5/0.165 = 57.6$ horizonten. (b) $180/57.6 = \qty{3.1}{m}$ kalksteen vandaag, $338/57.6 =
\qty{5.9}{m}$ vers sediment. (c) Hoogstens \qty{1.5}{m}: $180/1.5 = 120$ intervallen, 121
monsters. (d) Onderbrekingen in de afzetting, lagen waarin ammonieten ontbreken, en horizonten die tot
enkele centimeters zijn gecondenseerd waar het sediment traag aankwam.
\end{solution}

\begin{solution}{pb:b2:palaeoenvironments:1}
\textbf{1.} $866 \sin \ang{12} = 866 \times 0.2079 = \qty{180.0}{m}$.
\textbf{2.} $481 \times 0.2079 = \qty{100.0}{m}$.
\textbf{3.} De zandsteen, op grond van superpositie: de flanken hellen maar \ang{12}, dus de lagen zijn niet
omgekiept en de onderste is de oudste.
\textbf{4.} $866 - 766 = \qty{100}{m}$ en $1046 - 866 = \qty{180}{m}$.
\textbf{5.} Daar minstens \qty{766}{m}; zijn oorspronkelijke top is weggeërodeerd, dus zijn volledige dikte is
onbekend.
\textbf{6.} 0.70 aan het oppervlak; $0.70\,e^{-766/1408} = 0.406$ op \qty{766}{m}.
\textbf{7.} Het zijn stenohaliene groepen: het water had het zoutgehalte van de open zee.
\textbf{8.} Een \omterm{def:b2:palaeoenvironments:ichnofacies}{Cruziana-ichnofacies}: de offshorezone, tussen de golfbasis bij mooi weer en de
stormgolfbasis.
\textbf{9.} Stormen: hun golven bereikten de bodem, braken de schelpen en zetten ze af in lagen
met erosieve bases.
\textbf{10.} Het bodemwater was arm aan zuurstof: weinig gravers, laminae bewaard, pyriet gevormd.
\textbf{11.} De \omterm{def:b2:palaeoenvironments:ichnofacies}{Zoophycos-ichnofacies}, in het Mesozoïcum het lagere plat of de helling, onder de
stormgolfbasis; af en toe genoeg zuurstof voor gravers.
\textbf{12.} Nee: ammonieten zwommen in het water erboven; hun schelpen zonken of dreven na de dood naar de bodem
en vormen een \omterm{def:b2:palaeoenvironments:assemblages}{doodsassociatie}.
\textbf{13.} $184.2 - 174.7 = \qty{9.5}{Myr}$.
\textbf{14.} $180/9.5 = \qty{18.9}{m}$ per miljoen jaar.
\textbf{15.} $(2710 - 2573)/(2710 - 1000) = 0.080$.
\textbf{16.} $180 \times (1 - 0.080)/(1 - 0.51) = \qty{338}{m}$.
\textbf{17.} $338/9.5 = \qty{35.6}{m}$ per miljoen jaar.
\textbf{18.} $180 \times 0.92/9.5 = \qty{17.4}{m}$ per miljoen jaar.
\textbf{19.} De schalie gaf 11.3 en \qty{4.4}{m} per miljoen jaar: de korrelsnelheid daalde
ongeveer vier keer. Het carbonaat van de zandbank en het plat hield op met aankomen terwijl de bodem
dieper werd, en alleen een trage regen van modder bereikte hem.
\textbf{20.} Ze rukte op: rivieren, dan een zandbank, een offshorebodem, een bodem onder de
stormgolfbasis, dan de open zee van het krijt. De kolom legt een transgressie vast.
\textbf{21.} $57 - 41 = \qty{16}{m}$ tijdens de kalksteen, $86 - 57 = \qty{29}{m}$ tijdens de schalie,
\qty{45}{m} in totaal.
\textbf{22.} $(60 - 5) + 338 + 148.6 = \qty{542}{m}$, in $184.2 - 161.5 = \qty{22.7}{Myr}$:
\qty{23.9}{m} per miljoen jaar.
\textbf{23.} $45/542 = 0.083$: ongeveer \qty{8}{\percent}, en minder als de cyclusgrafiek
de amplitude overdrijft. De rest kwam van de daling van de bodem.
\textbf{24.} De kalkmodder van de kalksteen van de Varden Hills hoopte zich bij de afzetting op met ongeveer
\textbf{\qty[text-series-to-math=true]{36}{m}} per miljoen jaar.
\end{solution}
